\SourceLeaf{60}
$(1+n\cos\varphi)^{\nu-1}$ herleiten kann. So lässt sich von $(1+n\cos\varphi)^{-1}$ zu höheren negativen Potenzen übergehen, wobei man für die einzelnen Koeffizienten keine unendlichen Reihen benötigt.

$\alpha$) Um das für $\nu=-1$ geltende Gesetz des Anfangskoeffizienten $A$ zu untersuchen, setzen wir im allgemeinen Ausdruck für sein $h^{\text{ti}}$ Glied $\nu=-1$ ein. Wegen
\[
{}^{-1}\mathfrak{B}^{(2h)} = 1, \qquad \text{jenes Glied} \qquad = \frac{1}{2^{2h}}\,{}^{2h}\mathfrak{B}^{(h)}\,n^{2h}.
\]
Es ist aber
\[
\frac{1}{2^{2h}}\,{}^{2h}\mathfrak{B}^{(h)} = \frac{1}{2^{2h}}\,\frac{2h.(2h-1).\,.\,.\,.(h+1)}{1.2.\,.\,.\,. h}
\]
\[
= \frac{1}{2^{2h}} . \frac{2h.(2h-1).\,.\,.\,. 1}{1.2.\,.\,.\,.\,. h} . \frac{1}{h.(h-1).\,.\,. 1}
\]
\[
= \frac{1}{2^{2h}}\,\frac{(2h-1)(2h-3).\,.\,.\,. 1}{1.2.\,.\,. h} . \frac{2h.(2h-2).\,.\,. 2}{h.(h-1).\,.\,.\,.\,. 1}
\]
\[
= \frac{1}{2^{h}} . \frac{(2h-1)(2h-3).\,.\,. 1}{1.2.\,.\,.\,. h},
\]
und durch diese Umformungen wird jenes $h^{\text{tus}}$ Glied zu
\[
= \frac{1}{2^{h}}\,\frac{(2h-1).(2h-3).\,.\,. 1}{1.2.\,.\,.\,. h} .  n^{2h}.
\]
Da dieser Ausdruck offenbar in der Reihenentwicklung der Größe $(1-nn)^{-\frac{1}{2}}$ das $h^{\text{tus}}$ Glied ist, ergibt sich für $A$ selbst im Fall $\nu=-1$
\[
= (1-nn)^{-\frac{1}{2}}
\]
\par\noindent\hspace*{\parindent}$\beta$) Der erste Koeffizient nach dem Anfangskoeffizienten, $A^{(1)}$, wird durch die zuvor bewiesene Beziehung (7, $\delta$) bestimmt
\[
A^{(1)} = \frac{2.(\nu+2)}{n}\int An\,\mathrm{d}n - 2n.A,
\]
die hier wegen $\nu=-1$ übergeht in
\[
A^{(1)} = \frac{2}{n}\int An\,\mathrm{d}n - 2nA = \frac{2}{n}\int(1-nn)^{-\frac{1}{2}}n\,\mathrm{d}n - 2n(1-nn)^{-\frac{1}{2}}.
\]
Führt man die Integration so aus, dass das Integral mit $n$ verschwindet, so erhält man
\[
A^{(1)} = -\frac{2}{n}(1-nn)^{\frac{1}{2}} + \frac{2}{n} - 2n(1-nn)^{-\frac{1}{2}},
\]
oder kürzer
\[
A^{(1)} = \frac{2}{n}\bigl(1-(1-nn)^{-\frac{1}{2}}\bigr).
\]

